% A proof of Graffiti.pc Conjecture 143 (WOWII): t(G)*delta'(G) >= girth(G)+1.
\documentclass[11pt]{article}
\pdfoutput=1
\usepackage[margin=1.1in]{geometry}
\usepackage{amsmath,amssymb,amsthm}
\usepackage[colorlinks=true,linkcolor=blue,citecolor=blue,urlcolor=blue]{hyperref}
\numberwithin{equation}{section}

\newtheorem{theorem}{Theorem}[section]
\newtheorem{lemma}[theorem]{Lemma}
\newtheorem{proposition}[theorem]{Proposition}
\theoremstyle{remark}
\newtheorem{remark}[theorem]{Remark}

\newcommand{\tr}{\operatorname{t}}
\newcommand{\gr}{\operatorname{g}}
\newcommand{\dd}{\delta'}

\title{Largest induced trees, girth, and the second-smallest degree:\\
a proof of Graffiti.pc's Conjecture 143}
\author{Alper Ferudun\\\texttt{alper@mercurycodelab.com}}
\date{July 2026}

\begin{document}
\maketitle

\begin{abstract}
For a finite simple graph $G$, let $\tr(G)$ denote the largest number of
vertices inducing a tree in $G$, let $\gr(G)$ denote the girth of $G$, and let
$\dd(G)$ denote the second-smallest entry of the degree sequence of $G$,
counted with multiplicity.  Conjecture~143 of the conjecture-making program
Graffiti.pc (DeLaVi\~na, \emph{Written on the Wall II}, 2005) asserts that
every finite connected graph $G$ that is not a tree satisfies
$\tr(G)\ge(\gr(G)+1)/\dd(G)$.  We prove the conjecture.  The main step is an
elementary global argument showing that a connected graph containing a cycle
and at least two vertices of degree one has an induced tree on at least
$\gr(G)+1$ vertices.  The bound is sharp for every girth.  The proof has been
formalized and machine-checked in Lean~4 against the statement of the
conjecture in the Google DeepMind \emph{Formal Conjectures} repository.
\end{abstract}

\section{Introduction}

Graffiti.pc, developed by DeLaVi\~na~\cite{DeL01,DeLhistory} as a successor of
Fajtlowicz's program Graffiti~\cite{Faj88}, generates conjectural inequalities
between graph invariants.  Its conjectures are collected in the lists
\emph{Written on the Wall II}~\cite{WOWII}; a curated register with commentary
is maintained by West~\cite{WestReg}.  Many of these conjectures concern the
invariant $\tr(G)$, the largest order of an induced subgraph of $G$ that is a
tree, an invariant whose systematic study goes back to Erd\H{o}s, Saks, and
S\'os~\cite{ESS86}.

Throughout, $G$ is a finite simple graph.  We write $\gr(G)$ for the girth of
$G$ (the length of a shortest cycle), and $\dd(G)$ for the \emph{second-smallest
degree} of $G$: the second entry, with multiplicity, of the degree sequence of
$G$ sorted in nondecreasing order.  Thus $\dd(G)$ equals the minimum degree
$\delta(G)$ whenever at least two vertices attain the minimum degree.  A
\emph{leaf} is a vertex of degree one.

Conjecture~143 of Written on the Wall II, dated 2005, reads as
follows~\cite{WestReg}: \emph{if $G$ is connected and not a tree, then
$\tr(G)\ge(\gr(G)+1)/\dd(G)$.}  At the time of writing, the conjecture is
listed as open in West's register~\cite{WestReg} and is stated as a
research-open item in the Google DeepMind \emph{Formal Conjectures}
repository~\cite{FC,FC143}.  We prove it in the following equivalent,
denominator-free form.

\begin{theorem}\label{thm:main}
Let $G$ be a finite simple connected graph that is not a tree.  Then
\[
  \tr(G)\,\dd(G)\;\ge\;\gr(G)+1 .
\]
\end{theorem}

The proof splits on $\dd(G)$.  When $\dd(G)\ge2$, deleting one vertex of a
shortest (hence chordless) cycle leaves an induced path on $\gr(G)-1$ vertices,
and the inequality follows from $\gr(G)\ge3$ by arithmetic.  The substance of
the theorem is the case $\dd(G)=1$, in which $G$ has at least two leaves; here
we prove the following lemma, which may be of independent interest.

\begin{lemma}[Two-leaf lemma]\label{lem:twoleaf-intro}
Let $G$ be a finite simple connected graph that contains a cycle and has at
least two vertices of degree one.  Then $G$ has an induced tree on at least
$\gr(G)+1$ vertices.
\end{lemma}

The lemma is proved by a maximality argument: among the induced trees
containing two prescribed leaves, a maximum-order one admits an outside vertex
with two neighbours on the tree, which closes a cycle avoiding both leaves.
Both the lemma and the theorem are sharp for every girth
(Section~\ref{sec:sharp}).

\paragraph{Related work.}
The closest antecedent we are aware of is the induced-\emph{forest} bound of
DeLaVi\~na and Waller~\cite{DW04} (see also the discussion in
Hertz--Marcotte--Schindl~\cite{HMS14}): if $f(G)$ denotes the largest order of
an induced forest and $f_1(G)$ the number of leaves of a connected graph $G$
with a cycle, then $f(G)\ge\gr(G)+f_1(G)-1$.  With two leaves this produces an
induced forest on $\gr(G)+1$ vertices, but the forest need not be connected,
and the general forest-to-tree transfer of~\cite[Theorem~2.3]{HMS14} is too
lossy to recover Lemma~\ref{lem:twoleaf-intro} from it.  The foundational
paper~\cite{ESS86} bounds $\tr(G)$ in terms of order, size, radius,
independence and clique numbers, but contains no bound involving girth or the
second-smallest degree.

\paragraph{A remark on priority.}
Conjecture~143 has, to our knowledge, no published proof: it is carried as
open by the current registers~\cite{WestReg,FC143}, it does not appear on the
resolved-conjecture lists accompanying~\cite{WOWII}, and formula and citation
searches around~\cite{ESS86,DW04,HMS14} located no proof of the theorem or of
Lemma~\ref{lem:twoleaf-intro}.  While this work was being completed, a pull
request against the Formal Conjectures repository (\#4442, filed July~16,
2026, one day before the present proof was found) appeared, announcing a
machine-assisted resolution of the same conjecture via an externally hosted
Lean development.  That claim is unreviewed and unmerged at the time of
writing; we have, however, compiled the linked development against the
repository's pinned toolchain and confirmed that it does prove the
repository statement, so it constitutes an earlier concurrent resolution.
The present proof and formalization were obtained independently of it, and
this note appears to be the first human-readable exposition of a proof.
Because the argument is short and close in spirit to the 2004 forest bound,
we also do not claim that it could not have been observed before.

\paragraph{Verification.}
Beyond the human-readable proof below, Theorem~\ref{thm:main} has been
formalized and machine-checked in Lean~4~\cite{Lean4} on top of
Mathlib~\cite{mathlib}, against the pre-existing formal statement of
Conjecture~143 in the Formal Conjectures repository~\cite{FC143}; see
Section~\ref{sec:formal}.  As an independent falsification test, the theorem,
the lemma, and the constrained form of the maximality argument were verified
exhaustively (by two independent exact programs) on all $1252$ unlabeled
graphs on one to seven vertices, with no violation; the computation plays no
role in the proof.

\section{Notation}

All graphs are finite and simple.  For $S\subseteq V(G)$ we write $G[S]$ for
the induced subgraph.  A set $S$ \emph{induces a tree} if $G[S]$ is connected
and acyclic, and
\[
  \tr(G)\;=\;\max\{\,|S| : G[S]\ \text{is a tree}\,\}.
\]
A path or cycle \emph{in} $G$ is always a subgraph; a path $P$ is
\emph{induced} if $G[V(P)]=P$.  The girth $\gr(G)$ of a graph containing a
cycle is the minimum length of a cycle of $G$; note $\gr(G)\ge3$.  The degree
sequence of $G$ is the multiset of vertex degrees sorted in nondecreasing
order $d_1\le d_2\le\cdots\le d_n$, and $\dd(G)=d_2$.  This with-multiplicity
reading of ``second-smallest degree'' is the one fixed by the Graffiti.pc
definition list (Written on the Wall II, definition entry~65)~\cite{WOWII},
and it is the reading formalized in~\cite{FC143}.

We use two elementary facts.  First, in a connected graph on at least two
vertices every degree is positive; hence if $\dd(G)=d_2=1$ then $d_1=1$ as
well, so $G$ has at least two leaves.  Second, a shortest cycle $C$ of $G$ is
chordless, i.e.\ $G[V(C)]=C$: a chord splits $C$ into two cycles, each shorter
than $C$.

\section{The two-leaf lemma}

\begin{lemma}[Lemma~\ref{lem:twoleaf-intro}, strengthened]\label{lem:twoleaf}
Let $G$ be a finite simple connected graph that contains a cycle and has two
distinct vertices $x,y$ of degree one.  Then $G$ has an induced tree on at
least $\gr(G)+1$ vertices; in fact some induced tree of order at least
$\gr(G)+1$ contains both $x$ and $y$.
\end{lemma}

\begin{proof}
A shortest $x$--$y$ path in $G$ is induced (a chord would shorten it), and it
is a tree containing $x$ and $y$.  Hence the family of vertex sets
\[
  \mathcal{T}\;=\;\{\,S\subseteq V(G)\ :\ x,y\in S,\ G[S]\ \text{is a tree}\,\}
\]
is nonempty, and since $G$ is finite we may choose $S\in\mathcal{T}$ of
maximum cardinality.  Write $T=G[S]$.

The set $S$ is a proper subset of $V(G)$: if $S=V(G)$, then $G=T$ would be a
tree, contrary to the assumption that $G$ contains a cycle.  Since $G$ is
connected and $S$ is nonempty and proper, some edge of $G$ joins $S$ to its
complement; let $z\notin S$ be a vertex with a neighbour in $S$.

We claim $z$ has at least two neighbours in $S$.  Otherwise it has exactly
one, say $a$, and then $G[S\cup\{z\}]$ is the tree $T$ with the single pendant
vertex $z$ attached at $a$: it is connected, and acyclic because every cycle
of $G[S\cup\{z\}]$ through $z$ would need two distinct neighbours of $z$ in
$S$, while a cycle avoiding $z$ would lie in the tree $T$.  Thus
$S\cup\{z\}\in\mathcal{T}$ has larger cardinality than $S$, contradicting
maximality.

Choose distinct neighbours $a,b\in S$ of $z$, and let $P$ be the unique
$a$--$b$ path in the tree $T$.  The edges of $P$ together with $za$ and $zb$
form a (not necessarily induced) cycle of $G$ of length $|V(P)|+1$; possible
further edges from $z$ to $V(P)$ are irrelevant to its existence.  Hence
\begin{equation}\label{eq:girth}
  \gr(G)\;\le\;|V(P)|+1 .
\end{equation}

Finally, neither $x$ nor $y$ lies on $P$.  Indeed, every internal vertex of
$P$ has two distinct neighbours on $P$, and each endpoint ($a$ or $b$) has one
neighbour on $P$ and the additional neighbour $z\notin S$; so every vertex of
$P$ has degree at least two in $G$, whereas $\deg_G(x)=\deg_G(y)=1$.  Since
$V(P)\subseteq S$ and $x,y\in S$, the sets $V(P)$ and $\{x,y\}$ are disjoint
subsets of $S$, so by~\eqref{eq:girth}
\[
  \tr(G)\;\ge\;|S|\;\ge\;|V(P)|+2\;\ge\;\gr(G)+1 . \qedhere
\]
\end{proof}

\section{Proof of Theorem~\ref{thm:main}}

\begin{proof}[Proof of Theorem~\ref{thm:main}]
Let $G$ be connected and not a tree; then $G$ contains a cycle, so
$\gr(G)\ge3$, and $|V(G)|\ge3$, so all degrees are positive and $\dd(G)\ge1$.

\emph{Case $\dd(G)\ge2$.}  Let $C$ be a shortest cycle of $G$; as noted, $C$
is chordless.  Deleting one vertex of $C$ leaves an induced path on
$\gr(G)-1$ vertices, which is an induced tree, so $\tr(G)\ge\gr(G)-1$.  Hence
\[
  \tr(G)\,\dd(G)\;\ge\;2\,(\gr(G)-1)\;\ge\;\gr(G)+1,
\]
the last inequality being equivalent to $\gr(G)\ge3$.

\emph{Case $\dd(G)=1$.}  Since all degrees are positive and the second entry
of the sorted degree sequence equals one, the first entry equals one as well,
so $G$ has two distinct leaves.  Lemma~\ref{lem:twoleaf} gives
$\tr(G)\ge\gr(G)+1$, and multiplying by $\dd(G)=1$ finishes the proof.
\end{proof}

\begin{remark}
The formal statement of Conjecture~143 in~\cite{FC143} quantifies over
connected graphs on a finite vertex type with at least two vertices and
positive $\dd$, without excluding trees, and uses the Mathlib convention that
an acyclic graph has girth $0$.  That extension is immediate: for a connected
tree $G$ on $n\ge2$ vertices, the whole vertex set induces a tree, so
$\tr(G)\,\dd(G)\ge n\ge2>1=\gr(G)+1$.  (Already $\tr(G)\ge1$ suffices.)
\end{remark}

\section{Sharpness}\label{sec:sharp}

\begin{proposition}\label{prop:sharp}
For every $g\ge3$ there is a connected non-tree graph $G$ with $\gr(G)=g$,
$\dd(G)=1$, and $\tr(G)=g+1$.  Moreover, for every $g\ge3$ there is a
connected non-tree graph $H$ with $\gr(H)=g$, $\dd(H)=2$, and
$\tr(H)=g-1$, so the case split of the proof is tight as well.
\end{proposition}

\begin{proof}
For $G$, take a cycle $C_g$ and attach two pendant vertices (to arbitrary,
not necessarily distinct, cycle vertices).  Then $\gr(G)=g$ and $\dd(G)=1$.
At most two cycle vertices support the pendants; deleting a cycle vertex
supporting neither pendant leaves an induced subgraph on $g+1$ vertices that
is connected and acyclic, so $\tr(G)\ge g+1$.  No induced tree has $g+2$
vertices, since the only induced subgraph on all $g+2$ vertices is $G$
itself, which contains a cycle.  Hence $\tr(G)=g+1$, attaining equality in
Theorem~\ref{thm:main} and Lemma~\ref{lem:twoleaf}.

For $H$, take the cycle $C_g$ itself: $\dd(H)=2$, and the largest induced
trees are the paths obtained by deleting one vertex, of order $g-1$.  At
$g=3$ this attains equality in Theorem~\ref{thm:main},
$\tr(H)\,\dd(H)=2\cdot2=4=\gr(H)+1$.
\end{proof}

\section{Formalization}\label{sec:formal}

The Google DeepMind \emph{Formal Conjectures} project~\cite{FC} maintains
Lean~4 formalizations of open conjectures; Conjecture~143 was added in June
2026 as the statement \texttt{conjecture143} in the file
\texttt{FormalConjectures/WrittenOnTheWallII/GraphConjecture143.lean},
marked \texttt{research open}~\cite{FC143}.  The statement is the
denominator-free real-valued inequality
$\gr(G)+1\le\tr(G)\cdot\dd(G)$ for connected graphs on a finite
vertex type with at least two vertices, under the hypothesis $\dd(G)>0$, with
$\tr$ and $\dd$ given by the repository definitions
\texttt{largestInducedTreeSize} and \texttt{secondSmallestDegree} and with
Mathlib's $\mathbb{N}$-valued girth (which is $0$ on acyclic graphs).

\begin{sloppypar}
We have produced a complete, \texttt{sorry}-free Lean~4 proof of exactly this
statement, machine-checked with Lean toolchain v4.27.0 against current
Mathlib~\cite{mathlib}.  The development follows the paper proof, with the
supporting results proved as reusable graph-theoretic API
(\texttt{SimpleGraph} namespace): the one-vertex extension of an induced tree
along a unique neighbour (\texttt{IsTree.induce\_insert\_of\_unique\_adj});
induced trees from geodesics
(\texttt{Walk.induce\_support\_isTree\_of\_length\_eq\_dist}) and the induced
path obtained from a shortest cycle
(\texttt{girth\_sub\_one\_le\_largestInducedTreeSize}); the existence of a
maximum-order induced tree containing two prescribed vertices
(\texttt{exists\_maximum\_induced\_tree\_containing}); the boundary-vertex and
two-neighbour maximality arguments
(\texttt{Connected.exists\_adj\_finset\_compl},
\texttt{exists\_two\_adj\_of\_maximum\_induced\_tree\_containing}); the
cycle-closing certificate with the counting steps of Lemma~\ref{lem:twoleaf}
(\texttt{IsTree.girth\_add\_one\_le\_card\_of\_two\_leaves\_of\_two\_adj},
\texttt{girth\_add\_one\_le\_largestInducedTreeSize\_of\_two\_leaves}); and the
extraction of two leaves from $\dd(G)=1$
(\texttt{exists\_distinct\_degree\_one\_of\_secondSmallestDegree\_eq\_one}).
The main theorem is then assembled by the case split of
Theorem~\ref{thm:main} in under thirty lines.  The proof uses no
\texttt{native\_decide} and no additional axioms: auditing
\texttt{\#print axioms conjecture143} reports only \texttt{propext},
\texttt{Classical.choice}, and \texttt{Quot.sound}.  The full proof was
additionally re-derived and compiled in a second, independently written
assembly as a cross-check.  The source is available in the author's fork of
the repository and is included as ancillary files with this submission; a
pull request updating the status of Conjecture~143 is being coordinated with
the repository maintainers.
\end{sloppypar}

\paragraph{Acknowledgements.}
The author thanks the maintainers of the Formal Conjectures repository for
the formal statement, and Ermelinda DeLaVi\~na and Douglas B.\ West for
maintaining the Graffiti.pc conjecture lists.

\begin{thebibliography}{10}

\bibitem{DeL01}
E.~DeLaVi\~na,
\emph{Graffiti.pc},
Graph Theory Notes of New York XLII:3 (2002), 26--30.

\bibitem{DeLhistory}
E.~DeLaVi\~na,
\emph{Some history of the development of Graffiti},
in: Graphs and Discovery, DIMACS Ser.\ Discrete Math.\ Theoret.\ Comput.\
Sci.\ 69, AMS, 2005, 81--118.

\bibitem{WOWII}
E.~DeLaVi\~na,
\emph{Written on the Wall II: Conjectures of Graffiti.pc},
web list, University of Houston--Downtown,
\url{http://cms.dt.uh.edu/faculty/delavinae/research/wowII/} (retrieved July
2026).

\bibitem{WestReg}
D.~B. West,
\emph{Some conjectures of Graffiti.pc},
web register,
\url{https://dwest.web.illinois.edu/regs/graffiti.html} (retrieved July
2026).

\bibitem{Faj88}
S.~Fajtlowicz,
\emph{On conjectures of Graffiti},
Discrete Math.\ 72 (1988), 113--118.

\bibitem{ESS86}
P.~Erd\H{o}s, M.~Saks, and V.~T. S\'os,
\emph{Maximum induced trees in graphs},
J.~Combin.\ Theory Ser.\ B 41 (1986), 61--79.

\bibitem{DW04}
E.~DeLaVi\~na and B.~Waller,
\emph{On some conjectures of Graffiti.pc on the maximum order of induced
subgraphs},
Congr.\ Numer.\ 166 (2004), 11--32.

\bibitem{HMS14}
A.~Hertz, O.~Marcotte, and D.~Schindl,
\emph{On the maximum orders of an induced forest, an induced tree, and a
stable set},
Yugosl.\ J.\ Oper.\ Res.\ 24 (2014), no.~2, 199--215.
\href{https://doi.org/10.2298/YJOR130402037H}{doi:10.2298/YJOR130402037H}.

\bibitem{FC}
Google DeepMind,
\emph{Formal Conjectures},
GitHub repository,
\url{https://github.com/google-deepmind/formal-conjectures}.

\bibitem{FC143}
Formal Conjectures,
\emph{GraphConjecture143.lean},
\url{https://github.com/google-deepmind/formal-conjectures/blob/main/FormalConjectures/WrittenOnTheWallII/GraphConjecture143.lean}
(retrieved July 2026).

\bibitem{Lean4}
L.~de~Moura and S.~Ullrich,
\emph{The Lean 4 theorem prover and programming language},
in: Automated Deduction -- CADE 28, LNCS 12699, Springer, 2021, 625--635.

\bibitem{mathlib}
The mathlib Community,
\emph{The Lean mathematical library},
in: Proc.\ 9th ACM SIGPLAN Int.\ Conf.\ Certified Programs and Proofs
(CPP 2020), ACM, 2020, 367--381.

\end{thebibliography}

\end{document}
